% figA7_talignment.tex — Figure A7: T-alignment heatmap + ERM/PMH theory contrast
% % figA7_talignment.tex — Figure A7: T-alignment heatmap + ERM/PMH theory contrast
% % figA7_talignment.tex — Figure A7: T-alignment heatmap + ERM/PMH theory contrast
% % figA7_talignment.tex — Figure A7: T-alignment heatmap + ERM/PMH theory contrast
% \input{figures/figA7_talignment}
% Requires: colortbl, booktabs

\begin{figure*}[t]
\centering
\begin{subfigure}[b]{0.50\textwidth}
\centering
\small
\setlength{\tabcolsep}{5pt}
\renewcommand{\arraystretch}{1.18}
\begin{tabular}{lcccc}
\toprule
$\sigma_\mathrm{train}$ $\backslash$ $\sigma_\mathrm{eval}$
  & 0.05 & 0.10 & 0.15 & 0.20 \\
\midrule
0.05          & \cellcolor{colPMH!55}\textbf{79.36} & 73.24 & 59.97 & 46.08 \\
0.08          & 78.90 & \cellcolor{colPMH!48}\textbf{75.08} & 61.63 & 46.98 \\
0.10          & 78.92 & 75.68 & \cellcolor{colPMH!52}\textbf{68.16} & 54.45 \\
0.12 (def.)   & 78.71 & 75.27 & 70.63 & \cellcolor{colPMH!50}\textbf{59.00} \\
0.15          & 78.19 & 73.69 & \cellcolor{colPMH!50}\textbf{71.37} & 64.64 \\
0.20          & 77.99 & 74.61 & 71.16 & \cellcolor{colPMH!50}\textbf{69.05} \\
\bottomrule
\end{tabular}
\vspace{4pt}\\
{\footnotesize Shaded/bold $=$ column best. Every column peaks on diagonal --- zero exceptions.}
\caption{T-alignment heatmap (Task~04, ViT).}
\end{subfigure}
\hfill
\begin{subfigure}[b]{0.46\textwidth}
\centering
\begin{tikzpicture}[
  every node/.style={font=\footnotesize},
  box/.style={draw, rounded corners=4pt, inner sep=5pt, line width=0.6pt,
              minimum width=3.9cm, align=center},
]
\node[box, draw=colRed!60, fill=colRed!5] (erm) at (0,0) {
  \textbf{ERM (stuck)}\\[2pt]
  \tiny $\min\,\mathbb{E}[\mathcal{L}_\mathrm{task}(f(x),y)]$\\
  \tiny No constraint on geometry\\[3pt]
  \tiny \textit{Effect:} ERM must encode $n$\\
  \tiny $\|J_{\phi,n}w_n\|/L > 0$\\[3pt]
  \tiny \textit{Outcome:} Bound unavoidable\\[2pt]
  \normalsize\textbf{TDI@0 $= 1.093$}
};
\node[box, draw=colGreen!60, fill=colGreen!5] (pmh) at (0,-3.6) {
  \textbf{PMH (fixed)}\\[2pt]
  \tiny $+\,\|\phi(x){-}\phi(x{+}\delta)\|^2$,
    $\delta{\sim}\mathcal{N}(0,\sigma^2 I)$\\
  \tiny Uniquely isotropic (Prop.~5)\\[3pt]
  \tiny \textit{Effect:} Full Jacobian Fro suppressed\\
  \tiny uniformly in all directions\\[3pt]
  \tiny \textit{Outcome:} Bound broken $\checkmark$\\[2pt]
  \normalsize\textbf{TDI@0 $= 0.904$ ($-17.3\%$)}
};
\draw[-{Latex[length=4pt]}, colPMH, line width=1.0pt]
  (erm.south) -- (pmh.north)
  node[midway, right, font=\tiny, color=colPMH]
    {add PMH term};
\end{tikzpicture}
\caption{Why ERM is stuck and how PMH escapes.}
\end{subfigure}
\caption{\textbf{T-alignment and the ERM/PMH theory contrast.}
\emph{Left:} Every column peaks on the diagonal: training at $\sigma_\mathrm{train}$ is
optimal when evaluated at the matching $\sigma_\mathrm{eval}$, zero exceptions across 24 cells.
Asymmetry is $17\times$ (accuracy table; cf.\ multi-scale appendix for a TDI-ratio viewpoint):
under-suppression costs far more than over-suppression.
\emph{Right:} ERM is geometrically stuck (Theorem~1); one additional term breaks the bound
via isotropic Gaussian perturbations (Proposition~5).}
\label{fig:a7}
\end{figure*}

% Requires: colortbl, booktabs

\begin{figure*}[t]
\centering
\begin{subfigure}[b]{0.50\textwidth}
\centering
\small
\setlength{\tabcolsep}{5pt}
\renewcommand{\arraystretch}{1.18}
\begin{tabular}{lcccc}
\toprule
$\sigma_\mathrm{train}$ $\backslash$ $\sigma_\mathrm{eval}$
  & 0.05 & 0.10 & 0.15 & 0.20 \\
\midrule
0.05          & \cellcolor{colPMH!55}\textbf{79.36} & 73.24 & 59.97 & 46.08 \\
0.08          & 78.90 & \cellcolor{colPMH!48}\textbf{75.08} & 61.63 & 46.98 \\
0.10          & 78.92 & 75.68 & \cellcolor{colPMH!52}\textbf{68.16} & 54.45 \\
0.12 (def.)   & 78.71 & 75.27 & 70.63 & \cellcolor{colPMH!50}\textbf{59.00} \\
0.15          & 78.19 & 73.69 & \cellcolor{colPMH!50}\textbf{71.37} & 64.64 \\
0.20          & 77.99 & 74.61 & 71.16 & \cellcolor{colPMH!50}\textbf{69.05} \\
\bottomrule
\end{tabular}
\vspace{4pt}\\
{\footnotesize Shaded/bold $=$ column best. Every column peaks on diagonal --- zero exceptions.}
\caption{T-alignment heatmap (Task~04, ViT).}
\end{subfigure}
\hfill
\begin{subfigure}[b]{0.46\textwidth}
\centering
\begin{tikzpicture}[
  every node/.style={font=\footnotesize},
  box/.style={draw, rounded corners=4pt, inner sep=5pt, line width=0.6pt,
              minimum width=3.9cm, align=center},
]
\node[box, draw=colRed!60, fill=colRed!5] (erm) at (0,0) {
  \textbf{ERM (stuck)}\\[2pt]
  \tiny $\min\,\mathbb{E}[\mathcal{L}_\mathrm{task}(f(x),y)]$\\
  \tiny No constraint on geometry\\[3pt]
  \tiny \textit{Effect:} ERM must encode $n$\\
  \tiny $\|J_{\phi,n}w_n\|/L > 0$\\[3pt]
  \tiny \textit{Outcome:} Bound unavoidable\\[2pt]
  \normalsize\textbf{TDI@0 $= 1.093$}
};
\node[box, draw=colGreen!60, fill=colGreen!5] (pmh) at (0,-3.6) {
  \textbf{PMH (fixed)}\\[2pt]
  \tiny $+\,\|\phi(x){-}\phi(x{+}\delta)\|^2$,
    $\delta{\sim}\mathcal{N}(0,\sigma^2 I)$\\
  \tiny Uniquely isotropic (Prop.~5)\\[3pt]
  \tiny \textit{Effect:} Full Jacobian Fro suppressed\\
  \tiny uniformly in all directions\\[3pt]
  \tiny \textit{Outcome:} Bound broken $\checkmark$\\[2pt]
  \normalsize\textbf{TDI@0 $= 0.904$ ($-17.3\%$)}
};
\draw[-{Latex[length=4pt]}, colPMH, line width=1.0pt]
  (erm.south) -- (pmh.north)
  node[midway, right, font=\tiny, color=colPMH]
    {add PMH term};
\end{tikzpicture}
\caption{Why ERM is stuck and how PMH escapes.}
\end{subfigure}
\caption{\textbf{T-alignment and the ERM/PMH theory contrast.}
\emph{Left:} Every column peaks on the diagonal: training at $\sigma_\mathrm{train}$ is
optimal when evaluated at the matching $\sigma_\mathrm{eval}$, zero exceptions across 24 cells.
Asymmetry is $17\times$ (accuracy table; cf.\ multi-scale appendix for a TDI-ratio viewpoint):
under-suppression costs far more than over-suppression.
\emph{Right:} ERM is geometrically stuck (Theorem~1); one additional term breaks the bound
via isotropic Gaussian perturbations (Proposition~5).}
\label{fig:a7}
\end{figure*}

% Requires: colortbl, booktabs

\begin{figure*}[t]
\centering
\begin{subfigure}[b]{0.50\textwidth}
\centering
\small
\setlength{\tabcolsep}{5pt}
\renewcommand{\arraystretch}{1.18}
\begin{tabular}{lcccc}
\toprule
$\sigma_\mathrm{train}$ $\backslash$ $\sigma_\mathrm{eval}$
  & 0.05 & 0.10 & 0.15 & 0.20 \\
\midrule
0.05          & \cellcolor{colPMH!55}\textbf{79.36} & 73.24 & 59.97 & 46.08 \\
0.08          & 78.90 & \cellcolor{colPMH!48}\textbf{75.08} & 61.63 & 46.98 \\
0.10          & 78.92 & 75.68 & \cellcolor{colPMH!52}\textbf{68.16} & 54.45 \\
0.12 (def.)   & 78.71 & 75.27 & 70.63 & \cellcolor{colPMH!50}\textbf{59.00} \\
0.15          & 78.19 & 73.69 & \cellcolor{colPMH!50}\textbf{71.37} & 64.64 \\
0.20          & 77.99 & 74.61 & 71.16 & \cellcolor{colPMH!50}\textbf{69.05} \\
\bottomrule
\end{tabular}
\vspace{4pt}\\
{\footnotesize Shaded/bold $=$ column best. Every column peaks on diagonal --- zero exceptions.}
\caption{T-alignment heatmap (Task~04, ViT).}
\end{subfigure}
\hfill
\begin{subfigure}[b]{0.46\textwidth}
\centering
\begin{tikzpicture}[
  every node/.style={font=\footnotesize},
  box/.style={draw, rounded corners=4pt, inner sep=5pt, line width=0.6pt,
              minimum width=3.9cm, align=center},
]
\node[box, draw=colRed!60, fill=colRed!5] (erm) at (0,0) {
  \textbf{ERM (stuck)}\\[2pt]
  \tiny $\min\,\mathbb{E}[\mathcal{L}_\mathrm{task}(f(x),y)]$\\
  \tiny No constraint on geometry\\[3pt]
  \tiny \textit{Effect:} ERM must encode $n$\\
  \tiny $\|J_{\phi,n}w_n\|/L > 0$\\[3pt]
  \tiny \textit{Outcome:} Bound unavoidable\\[2pt]
  \normalsize\textbf{TDI@0 $= 1.093$}
};
\node[box, draw=colGreen!60, fill=colGreen!5] (pmh) at (0,-3.6) {
  \textbf{PMH (fixed)}\\[2pt]
  \tiny $+\,\|\phi(x){-}\phi(x{+}\delta)\|^2$,
    $\delta{\sim}\mathcal{N}(0,\sigma^2 I)$\\
  \tiny Uniquely isotropic (Prop.~5)\\[3pt]
  \tiny \textit{Effect:} Full Jacobian Fro suppressed\\
  \tiny uniformly in all directions\\[3pt]
  \tiny \textit{Outcome:} Bound broken $\checkmark$\\[2pt]
  \normalsize\textbf{TDI@0 $= 0.904$ ($-17.3\%$)}
};
\draw[-{Latex[length=4pt]}, colPMH, line width=1.0pt]
  (erm.south) -- (pmh.north)
  node[midway, right, font=\tiny, color=colPMH]
    {add PMH term};
\end{tikzpicture}
\caption{Why ERM is stuck and how PMH escapes.}
\end{subfigure}
\caption{\textbf{T-alignment and the ERM/PMH theory contrast.}
\emph{Left:} Every column peaks on the diagonal: training at $\sigma_\mathrm{train}$ is
optimal when evaluated at the matching $\sigma_\mathrm{eval}$, zero exceptions across 24 cells.
Asymmetry is $17\times$ (accuracy table; cf.\ multi-scale appendix for a TDI-ratio viewpoint):
under-suppression costs far more than over-suppression.
\emph{Right:} ERM is geometrically stuck (Theorem~1); one additional term breaks the bound
via isotropic Gaussian perturbations (Proposition~5).}
\label{fig:a7}
\end{figure*}

% Requires: colortbl, booktabs

\begin{figure*}[t]
\centering
\begin{subfigure}[b]{0.50\textwidth}
\centering
\small
\setlength{\tabcolsep}{5pt}
\renewcommand{\arraystretch}{1.18}
\begin{tabular}{lcccc}
\toprule
$\sigma_\mathrm{train}$ $\backslash$ $\sigma_\mathrm{eval}$
  & 0.05 & 0.10 & 0.15 & 0.20 \\
\midrule
0.05          & \cellcolor{colPMH!55}\textbf{79.36} & 73.24 & 59.97 & 46.08 \\
0.08          & 78.90 & \cellcolor{colPMH!48}\textbf{75.08} & 61.63 & 46.98 \\
0.10          & 78.92 & 75.68 & \cellcolor{colPMH!52}\textbf{68.16} & 54.45 \\
0.12 (def.)   & 78.71 & 75.27 & 70.63 & \cellcolor{colPMH!50}\textbf{59.00} \\
0.15          & 78.19 & 73.69 & \cellcolor{colPMH!50}\textbf{71.37} & 64.64 \\
0.20          & 77.99 & 74.61 & 71.16 & \cellcolor{colPMH!50}\textbf{69.05} \\
\bottomrule
\end{tabular}
\vspace{4pt}\\
{\footnotesize Shaded/bold $=$ column best. Every column peaks on diagonal --- zero exceptions.}
\caption{T-alignment heatmap (Task~04, ViT).}
\end{subfigure}
\hfill
\begin{subfigure}[b]{0.46\textwidth}
\centering
\begin{tikzpicture}[
  every node/.style={font=\footnotesize},
  box/.style={draw, rounded corners=4pt, inner sep=5pt, line width=0.6pt,
              minimum width=3.9cm, align=center},
]
\node[box, draw=colRed!60, fill=colRed!5] (erm) at (0,0) {
  \textbf{ERM (stuck)}\\[2pt]
  \tiny $\min\,\mathbb{E}[\mathcal{L}_\mathrm{task}(f(x),y)]$\\
  \tiny No constraint on geometry\\[3pt]
  \tiny \textit{Effect:} ERM must encode $n$\\
  \tiny $\|J_{\phi,n}w_n\|/L > 0$\\[3pt]
  \tiny \textit{Outcome:} Bound unavoidable\\[2pt]
  \normalsize\textbf{TDI@0 $= 1.093$}
};
\node[box, draw=colGreen!60, fill=colGreen!5] (pmh) at (0,-3.6) {
  \textbf{PMH (fixed)}\\[2pt]
  \tiny $+\,\|\phi(x){-}\phi(x{+}\delta)\|^2$,
    $\delta{\sim}\mathcal{N}(0,\sigma^2 I)$\\
  \tiny Uniquely isotropic (Prop.~5)\\[3pt]
  \tiny \textit{Effect:} Full Jacobian Fro suppressed\\
  \tiny uniformly in all directions\\[3pt]
  \tiny \textit{Outcome:} Bound broken $\checkmark$\\[2pt]
  \normalsize\textbf{TDI@0 $= 0.904$ ($-17.3\%$)}
};
\draw[-{Latex[length=4pt]}, colPMH, line width=1.0pt]
  (erm.south) -- (pmh.north)
  node[midway, right, font=\tiny, color=colPMH]
    {add PMH term};
\end{tikzpicture}
\caption{Why ERM is stuck and how PMH escapes.}
\end{subfigure}
\caption{\textbf{T-alignment and the ERM/PMH theory contrast.}
\emph{Left:} Every column peaks on the diagonal: training at $\sigma_\mathrm{train}$ is
optimal when evaluated at the matching $\sigma_\mathrm{eval}$, zero exceptions across 24 cells.
Asymmetry is $17\times$ (accuracy table; cf.\ multi-scale appendix for a TDI-ratio viewpoint):
under-suppression costs far more than over-suppression.
\emph{Right:} ERM is geometrically stuck (Theorem~1); one additional term breaks the bound
via isotropic Gaussian perturbations (Proposition~5).}
\label{fig:a7}
\end{figure*}
